\glsresetall

\Gls{dm} is an invisible form of matter whose subatomic nature remains
unknown, motivating searches for physics beyond the \gls{sm}. The \gls{ldmx}
searches for sub-GeV dark matter through a missing-momentum signature in
electron fixed-target interactions. Invisible-particle production would
appear as an imbalance between the incoming electron momentum and the
reconstructed visible final state. Reliable reconstruction of beam-electron
activity is therefore essential
\cite{cirelli2026-darkmatter,akesson2018-ldmx,ldmx2025-overview}.

In experimental particle physics, \emph{reconstruction} uses detector
measurements to infer physically meaningful properties of the underlying
particle event. At increased beam intensity, several electrons can enter the same detector
readout window and produce overlapping showers in the \gls{ecal}. Their
detector activity can then become difficult to associate with individual
electrons. Rejecting all multi-electron events would discard part of the
delivered beam, while reliable reconstruction could make a larger fraction of
recorded events useful and improve the practical physics reach of the
experiment \cite{ldmx2025-overview}.

This thesis investigates modern supervised \gls{ml} methods for separating showers
in simulated \(2e\) and \(3e\) events. The principal task assigns each
selected \gls{ecal} hit to the electron responsible for the largest simulated
deposited-energy contribution. Each event is represented as an unordered,
variable-size collection of reconstructed \gls{ecal} hits and
\gls{tpad}-track observations from the \gls{ts}. This preserves the irregular
detector structure and allows both subdetectors to be processed jointly
without projecting the event onto a fixed image.

A GravNet \gls{gnn} and a Transformer encoder provide two fixed-multiplicity
baseline architectures. Their reconstruction quality, model size and observed
training behaviour are compared, and sensitivity to the \gls{tpad} input is
evaluated through ablation. An exploratory \gls{mlpf}-inspired model,
named the \gls{trace}, extends the task to a mixed \(2e/3e\) sample. It predicts electron-slot
validity, the set of electron contributors to each hit and their
deposited-energy fractions. Electron count and dominant hit origin are then
derived consistently from these outputs
\cite{qasim2019-gravnet,pata2021-mlpf,pata2026-mlpf}.

The study is an offline proof of concept based on simulated non-signal events.
It does not perform a dark-matter search, establish detector-data performance
or deploy an online trigger algorithm, because that is out of scope. A further contribution is the reusable
\texttt{ml\_ldmx} platform, which connects LDMX data preparation, supervised
target construction, model training and reproducible evaluation and future studies, and is intended to support future extensions of the reconstruction methods developed here \cite{ml_ldmx-github}.

\clearpage
\subsection{Aim and Study Strategy}

The work was organised around the following objectives:

\begin{itemize}
    \item \textbf{Investigate shower separation:} determine how reliably
    supervised models can assign \gls{ecal} hits to their originating electrons in
    simulated \(2e\) and \(3e\) pile-up events.

    \item \textbf{Build a reusable research platform:} develop the
    \texttt{ml\_ldmx} pipeline for extracting LDMX simulation data,
    constructing truth targets, training models and evaluating predictions.

    \item \textbf{Compare model architectures:} evaluate modern set-based neural networks choices GravNet and Transformer baselines using common inputs and targets, considering both
    reconstruction quality and practical computational properties.

    \item \textbf{Identify limiting conditions:} relate reconstruction
    performance to electron multiplicity, shower overlap, \gls{ecal} depth and model
    confidence.

    \item \textbf{Assess heterogeneous detector information:} test whether
    reconstructed trigger-pad observations provide useful context beyond the
    \gls{ecal} measurements.

    \item \textbf{Explore a broader reconstruction task:} extend the fixed electron-count hit assignment with the mixed-multiplicity \gls{trace} model, inspired by machine-learning-based particle flow at CMS, which predicts electron-slot validity, contributing-electron sets and energy fractions. Electron count is derived from the valid slots.
\end{itemize}

The intended outcome is an offline proof-of-concept and a documented basis for
continued reconstruction studies, rather than a production-grade trigger algorithm,
a complete detector reconstruction chain or a dark-matter sensitivity study.

\glsresetall